% Part: first-order-logic
% Chapter: beyond
% Section: introduction

\documentclass[../../../include/open-logic-section]{subfiles}

\begin{document}

\olfileid{fol}{byd}{int}

\olsection{Ringkesan}

Logika tataran kapisan dudu siji-sijine sistem logika
kang narik kawigaten: ana akeh pangembangan
lan variasi saka logika tataran kapisan. Sawijining
logika lumrahe ngemot panetepan formal sawijining
basa, lan biasane, sanajan ora mesthi,
sistem deduktif, uga biasane, sanajan ora mesthi,
semantik kang dikarepake. Nanging panganggone
tembung iki sacara teknis nuwuhake pitakon
kang cetha: apa gegayutane sistem-sistem logika
kang dudu logika tataran kapisan karo tembung "logika"
ing pangertèn intuisi utawa filsafat? Kabeh
sistem kang bakal dijlentrehake ing ngisor iki
dirancang kanggo nggambarake sawenehing wujud
panalaran; apa kita bisa nerangake apa kang
ndadekake sistem-sistem iku logis?

Ora ana jawaban kang gampang. Tembung "logika"
digunakake kanthi maneka cara lan ing maneka
konteks, lan gagasan iki, kaya dene gagasan
"bebener", wis dianalisis saka akeh sudut
pandang filsafat. Upamane, tujuan panalaran
logis bisa dianggep minangka nemtokake
pratelan endi kang mesthi bener, bener
secara a priori, bener tanpa gumantung
marang interpretasi term nonlogis, bener
amarga wangune, utawa bener amarga
konvensi basa; saben pamahaman iki
mbutuhake panjlentrehan kang akeh. Sanajan
sing digatekake mung jinis logika kang
digunakake ing matematika, isih durung
akeh kasapukan ngenani cakupane. Upamane,
ing \textit{Principia Mathematica}, Russell
lan Whitehead nyoba ngembangake matematika
adhedhasar logika, ngetutake tradhisi
{\em logisis} kang diwiwiti Frege. Sistem
logikane iku sawijine wujud logika tipe
luwih dhuwur kang mirip karo kang bakal
dijlentrehake ing ngisor iki. Pungkasane,
padha kudu ngenalake aksioma kang miturut
akeh ukuran ora katon murni logis
(mligine aksioma infinitas lan aksioma
reduksibilitas), nanging wong isih bisa
nganggep sawenehing wujud panalaran
orde luwih dhuwur minangka logis.
Kosok balene, Quine, kang ontologine
ora ngakoni "proposisi" minangka objek
kang sah kanggo dirembug, duwe panemu
yen logika orde kapindho lan orde
luwih dhuwur sejatiné wujud teori
himpunan kang nyamar dadi logika;
kanthi tembung liya, sistem kang
ngemot kuantifikasi tumrap predikat
dudu logika murni.

Saiki, luwih becik prakara filsafat kaya
mangkono disisihake dhisik kanggo wektu
liyane, lan sistem-sistem ing ngisor iki
dianggep minangka idealisasi formal
saka maneka jinis panalaran, logis
utawa ora.

\end{document}
